\documentclass[a4paper]{article}

% Paquetes básicos
% Español (México) con XeLaTeX: fontspec para fuentes Unicode, polyglossia
% para la división silábica y los nombres del documento en la variante mexicana.
\usepackage{fontspec}
\usepackage{polyglossia}
\setdefaultlanguage[variant=mexican]{spanish}
% Los nombres chinos van como «pinyin (original)»: xeCJK da a los caracteres
% Han su propia fuente; sin ella XeLaTeX los omite con un aviso.
% FandolSong no cubre algunos caracteres de nombres (U+73FA); WenQuanYi Zen Hei sí.
\usepackage[AutoFallBack=true]{xeCJK}
\setCJKmainfont{FandolSong-Regular.otf}
\setCJKfallbackfamilyfont{\CJKrmdefault}{WenQuanYi Zen Hei}
% polyglossia no traduce los nombres de listados.
\AtBeginDocument{\providecommand{\lstlistingname}{}\renewcommand{\lstlistingname}{Listado}%
  \providecommand{\lstlistlistingname}{}\renewcommand{\lstlistlistingname}{Índice de listados}}
\usepackage{amsmath, amssymb}
\usepackage{graphicx}
\usepackage[margin=2.5cm]{geometry}
\usepackage[most]{tcolorbox}
\usepackage{etoolbox}
\usepackage{listings}
\usepackage{hyperref}
\usepackage{booktabs}
\usepackage{subcaption}
\usepackage{float}

% Cajas de resaltado
\newtcolorbox{knowledgebox}[1]{
    enhanced, colback=blue!5!white, colframe=blue!75!black, colbacktitle=blue!75!black,
    coltitle=white, fonttitle=\bfseries, title={#1}, attach boxed title to top left={yshift=-2mm, xshift=2mm},
    boxrule=1pt, sharp corners
}

\newtcolorbox{importantbox}[1]{
    enhanced, colback=yellow!10!white, colframe=yellow!80!black, colbacktitle=yellow!80!black,
    coltitle=black, fonttitle=\bfseries, title={#1}, sharp corners
}

\newtcolorbox{warningbox}[1]{
    enhanced, colback=red!5!white, colframe=red!75!black, colbacktitle=red!75!black,
    coltitle=white, fonttitle=\bfseries, title={#1}, sharp corners
}

% Listados
\lstset{
    language=Python,
    basicstyle=\ttfamily\small,
    keywordstyle=\color{blue},
    stringstyle=\color{red!60!black},
    commentstyle=\color{green!60!black},
    breaklines=true,
    frame=single,
    numbers=left,
    numberstyle=\tiny\color{gray},
    captionpos=b,
    extendedchars=true
}

% Metadatos de la página inicial
\newcommand{\notetitle}{Modelos de lenguaje difusivos: desde MDLM hasta SoLM}
\newcommand{\noteauthors}{Elaboradas a partir de materiales públicos del curso}
\newcommand{\notedate}{Otoño de 2025}
\newcommand{\videochannel}{Subham Sahoo (Clase invitada)}
\newcommand{\videopublishdate}{2025}
\newcommand{\videoduration}{40 minutos}
\newcommand{\videourl}{https://www.youtube.com/watch?v=eA1EnwoMxf0}
\newcommand{\videocoverpath}{cover.jpg}
\newcommand{\repourl}{https://github.com/hqhq1025/ai-course-notes}

\begin{document}

````


\begin{titlepage}
\centering
{\Large Notas del curso\par}
\vspace{1.2cm}
{\huge\bfseries \notetitle\par}
\vspace{0.8cm}
{\large \noteauthors\par}
\vspace{0.3cm}
{\large \notedate\par}
\vspace{1.2cm}

\ifdefempty{\videocoverpath}{
{\small Al generar las notas, indique la ruta local de la portada del video.\par}
}{
\includegraphics[width=0.82\textwidth,height=0.45\textheight,keepaspectratio]{\videocoverpath}\par
}

\vfill
\begin{tcolorbox}[width=0.9\textwidth, colback=black!2!white, colframe=black!60, sharp corners]
\textbf{Autor o canal del video}: \videochannel\par
\textbf{Fecha de publicación}: \videopublishdate\par
\textbf{Duración del video}: \videoduration\par
\ifdefempty{\videourl}{}{
\textbf{Enlace del video}: \href{\videourl}{\nolinkurl{\videourl}}\par
}
\end{tcolorbox}
\end{titlepage}

\tableofcontents
\newpage

%% --- Inicio del contenido --- %%

\section{Introducción: ¿por qué se necesitan modelos de lenguaje difusos}

Los modelos autoregresivos (Autoregressive, AR) han sido durante mucho tiempo el paradigma predominante en el campo de la generación de lenguaje. Desde la serie GPT hasta LLaMA, estos modelos generan texto palabra por palabra mediante un proceso estricto de izquierda a derecha (left-to-right), logrando un enorme éxito. Sin embargo, los modelos AR presentan una limitación fundamental de velocidad: generar $N$ tokens requiere $N$ pasos de propagación hacia adelante, lo que impide la generación verdadera en paralelo.

\begin{importantbox}{Ventaja central de los modelos de lenguaje difusos}
Los modelos de lenguaje difusos (Diffusion Language Models) se liberan de la restricción de izquierda a derecha y pueden generar texto en cualquier orden; lo más importante es que permiten \textbf{generar múltiples palabras en paralelo}. Esto les otorga una ventaja significativa en términos de velocidad de inferencia.
\end{importantbox}

Desde 2025, los modelos de lenguaje difusos han generado un interés generalizado en la industria. Google lanzó Gemini Diffusion, ByteDance introdujo Seed Diffusion e Inception Labs publicó Mercury Coder. Estos modelos no solo ofrecen un rendimiento excepcional, sino que superan ampliamente a los modelos AR tradicionales en velocidad de inferencia. Medios comerciales principales como Forbes y Fortune también han comenzado a reportar sobre esta transformación tecnológica.

El conferencista invitado de esta sesión, Subham Sahoo (萨哈姆·萨胡), es doctor en la Universidad Cornell; su investigación doctoral se centró en los modelos de lenguaje difusos. Es el primer autor de trabajos importantes como MDLM (Masked Diffusion Language Models) y SoLM (Sorting-based Language Models). Esta charla abordará a fondo las ideas centrales y los detalles técnicos de ambos artículos.

\subsection{Línea de tiempo del desarrollo de los modelos de lenguaje difusos}

\begin{knowledgebox}{De 2022 a 2025: hitos clave}
\begin{enumerate}
    \item \textbf{2022}: Se propuso por primera vez el modelo de lenguaje difuso. Estos modelos podían generar palabras en un orden aleatorio, pero la calidad del texto quedaba muy atrás respecto a los modelos AR, recibiendo una reacción fría por parte de la comunidad.
    \item \textbf{Principios de 2024}: Lou et al. demostraron que los modelos de lenguaje difusos combinados con arquitecturas Transformer modernas pueden igualar la calidad de los modelos AR a nivel de GPT-2, al mismo tiempo que son significativamente más rápidos.
    \item \textbf{Mediados de 2024}: Se publicó MDLM, que rederivó y simplificó el algoritmo de entrenamiento de difusión discreto, logrando un rendimiento comparable al de los modelos AR en conjuntos de datos de texto de pequeña escala.
    \item \textbf{Principios de 2025}: El artículo Diffusion Duality propuso un algoritmo de destilación que permite a los modelos generar secuencias de 1000 palabras con solo 10 pasos, manteniendo una calidad casi sin pérdida.
    \item \textbf{2025}: Se publicó SoLM, integrando AR y difusión discreta, con soporte para caché KV, lo que logró un aumento de orden de magnitud en la velocidad de inferencia.
\end{enumerate}
\end{knowledgebox}

\subsection{Resumen de la sección}

Los modelos de lenguaje difusos representan una transformación importante en el paradigma de generación de texto. Desde la propuesta conceptual en 2022 hasta el despliegue industrial en 2025, este campo ha experimentado un desarrollo rápido. El marco MDLM es actualmente la base de varios modelos de lenguaje difusos a nivel industrial (como Seed Diffusion de ByteDance).

\newpage
\section{MDLM: Modelo de lenguaje con difusión enmascarada}

\subsection{Objetivos del diseño e ideas centrales}

El objetivo de MDLM (Masked Diffusion Language Model) es diseñar un modelo de lenguaje que soporte la \textbf{generación paralela}. A diferencia de los modelos AR, que generan palabra por palabra, MDLM parte de una secuencia completamente enmascarada y recupera progresivamente el texto completo.

Intuición del proceso de generación:
\begin{enumerate}
    \item Estado inicial: una secuencia de longitud $L$ con todos los elementos enmascarados $[\text{M}, \text{M}, \ldots, \text{M}]$
    \item Se introduce la secuencia en un Transformer bidireccional similar a BERT
    \item El modelo predice las palabras para todas las posiciones enmascaradas
    \item Se revelan selectivamente algunas predicciones y el resto se vuelve a enmascarar
    \item Se repiten los pasos 2--4 hasta que todas las posiciones estén reveladas
\end{enumerate}

\begin{importantbox}{Diferencia fundamental entre MDLM y AR}
\begin{itemize}
    \item \textbf{Modelos AR}: utilizan \textbf{atención causal}, donde cada token depende únicamente del contexto a su izquierda, generando uno por uno
    \item \textbf{MDLM}: utiliza \textbf{atención bidireccional}, donde cada token depende del contexto completo, permitiendo generar múltiples palabras simultáneamente
\end{itemize}
\end{importantbox}

\subsection{Proceso hacia adelante: ruido enmascarado}

El marco central de los modelos de difusión consiste en definir un \textbf{proceso de corrosión hacia adelante} conocido $q$ y luego aprender su \textbf{proceso inverso} $p_\theta$.

En MDLM, el proceso hacia adelante es la operación de \textbf{enmascaramiento}. Dado un texto limpio $\mathbf{x}$, se genera una secuencia con ruido mediante el enmascaramiento progresivo de las palabras.

\begin{importantbox}{Nivel de señal}
MDLM utiliza un nivel de señal $\alpha_t \in [0, 1]$ para controlar el grado de ruido:
\begin{itemize}
    \item $\alpha_t = 1$: oración completamente limpia $\mathbf{x}$ (tiempo $t = 0$)
    \item $\alpha_t = 0$: secuencia completamente enmascarada (tiempo $t = 1$)
\end{itemize}
El proceso desde la oración limpia hasta el estado de ruido intermedio $\mathbf{z}_t$ es muy simple: para cada palabra limpia se lanza una moneda con probabilidad $1 - \alpha_t$ de reemplazarla por un token enmascarado.
\end{importantbox}

Formalmente, el proceso hacia adelante se expresa como:

$$q(\mathbf{z}_t | \mathbf{x}) = \prod_{i=1}^{L} q(z_t^{(i)} | x^{(i)})$$

donde cada posición opera independientemente:

$$q(z_t^{(i)} | x^{(i)}) = \begin{cases} \alpha_t & \text{si } z_t^{(i)} = x^{(i)} \text{(se mantiene la palabra original)} \\ 1 - \alpha_t & \text{si } z_t^{(i)} = \text{M} \text{(se reemplaza por enmascaramiento)} \end{cases}$$

\begin{itemize}
    \item $\mathbf{x}$: secuencia de texto limpio
    \item $\mathbf{z}_t$: secuencia de ruido correspondiente al tiempo $t$
    \item $\alpha_t$: nivel de señal en el tiempo $t$
    \item $L$: longitud de la secuencia
    \item M: token especial de enmascaramiento
\end{itemize}

\subsection{Posterior inversa: función de desenmascaramiento}

\begin{importantbox}{Posterior inversa}
Esta es una de las fórmulas más importantes del curso. La posterior inversa $q(\mathbf{z}_s | \mathbf{z}_t, \mathbf{x})$ describe cómo pasar de un estado más ruidoso $\mathbf{z}_t$ a uno más limpio $\mathbf{z}_s$ (donde $s < t$), dado el secuencia con ruido $\mathbf{z}_t$ y el texto limpio $\mathbf{x}$.
\end{importantbox}

Para cada posición $i$, el comportamiento de la posterior inversa depende del estado actual de dicha posición:
\begin{itemize}
    \item \textbf{Token revelado} ($z_t^{(i)} \neq \text{M}$): permanece inalterado, es decir, $z_s^{(i)} = z_t^{(i)}$
    \item \textbf{Token enmascarado} ($z_t^{(i)} = \text{M}$): se desenmascara a la palabra limpia $x^{(i)}$ con probabilidad $\frac{\alpha_s - \alpha_t}{1 - \alpha_t}$; de lo contrario, permanece enmascarado
\end{itemize}

\begin{warningbox}{La posterior inversa por sí misma no se puede usar directamente para generar}
Obsérvese que $q(\mathbf{z}_s | \mathbf{z}_t, \mathbf{x})$ requiere conocer el texto limpio $\mathbf{x}$, pero durante la generación desconocemos exactamente cuál es $\mathbf{x}$. Por tanto, el desafío consiste en aprender una \textbf{red de desruido} $x_\theta(\mathbf{z}_t, t)$ que estime $\mathbf{x}$ y, posteriormente, sustituir ese estimado en la fórmula de la posterior inversa para obtener $p_\theta(\mathbf{z}_s | \mathbf{z}_t)$.
\end{warningbox}

\subsection{Proceso de generación $p_\theta$}

Durante la inferencia, utilizamos el Transformer desruido entrenado $x_\theta$ para aproximar el proceso inverso:

\begin{enumerate}
    \item Se introduce la secuencia con ruido $\mathbf{z}_t$ en un Transformer bidireccional
    \item El modelo produce una \textbf{distribución de categorías} sobre el vocabulario para cada posición enmascarada
    \item Se muestrea de dicha distribución y se rellenan todas las posiciones enmascaradas
    \item Se vuelven a enmascarar selectivamente algunas predicciones (dependiendo del paso temporal, se decide cuántas retener)
    \item Se repite el proceso anterior hasta que todas las posiciones estén reveladas
\end{enumerate}

\begin{knowledgebox}{Operación de copiado}
El Transformer de MDLM incluye un mecanismo de copia: para las posiciones que ya están limpias (no enmascaradas), el modelo copia directamente la entrada como salida sin modificaciones. Esto implica que la pérdida de entropía cruzada se calcula efectivamente \textbf{únicamente en las posiciones enmascaradas}, lo cual simplifica considerablemente el entrenamiento.
\end{knowledgebox}

\subsection{Resumen de la sección}

El diseño central de MDLM puede resumirse así: (1) el proceso hacia adelante convierte el texto en una secuencia totalmente enmascarada mediante enmascaramiento progresivo; (2) se aprende un Transformer bidireccional para estimar el texto limpio; (3) se utiliza la fórmula de la posterior inversa para desenmascarar progresivamente y generar el texto.

\newpage
\section{Entrenamiento de MDLM}

\subsection{Función de pérdida ELBO}

En los modelos AR, el objetivo del entrenamiento es maximizar directamente la verosimilitud logarítmica $\log p(\mathbf{x})$. En los modelos de difusión no es posible calcular con precisión la verosimilitud, por lo que se utiliza el \textbf{límite inferior de evidencia} (Evidence Lower Bound, ELBO) como objetivo del entrenamiento.

Una de las principales contribuciones de MDLM es haber derivado un ELBO en tiempo continuo \textbf{extremadamente conciso}:

$$\mathcal{L}_{\text{MDLM}} = \mathbb{E}_{t \sim \mathcal{U}(0,1), \, \mathbf{z}_t \sim q(\mathbf{z}_t|\mathbf{x})} \left[ \frac{\alpha_t'}{1 - \alpha_t} \sum_{i: z_t^{(i)} = \text{M}} -\log x_\theta^{(i)}(\mathbf{z}_t, t) \right]$$

Significado de cada término:
\begin{itemize}
    \item $\mathbb{E}_{t \sim \mathcal{U}(0,1)}$: Se muestra el paso de tiempo $t$ uniformemente de $[0, 1]$
    \item $\mathbf{z}_t \sim q(\mathbf{z}_t | \mathbf{x})$: Se añade ruido al texto limpio $\mathbf{x}$ para obtener $\mathbf{z}_t$
    \item $\frac{\alpha_t'}{1 - \alpha_t}$: \textbf{Término de ponderación}, es decir, la probabilidad de desmascaramiento (chance of unmasking)
    \item $-\log x_\theta^{(i)}(\mathbf{z}_t, t)$: \textbf{Pérdida de entropía cruzada}, calculada únicamente en las posiciones con máscara
\end{itemize}

\begin{importantbox}{Límite continuo}
La idea clave del artículo sobre MDLM es que, cuando el número de pasos de discretización $T \to \infty$, la cadena de Markov de tiempo discreto converge a una cadena de Markov de tiempo continuo, obteniendo así la función de pérdida concisa anterior. Esta simplificación facilita enormemente el entrenamiento y la evaluación.
\end{importantbox}

\subsection{Comparación entre el entrenamiento AR y MDLM}

Los flujos de entrenamiento de los modelos AR y MDLM son muy similares; solo existen tres diferencias clave:

\begin{table}[H]
\centering
\caption{Comparación del entrenamiento de modelos AR y MDLM}
\begin{tabular}{@{}lll@{}}
\toprule
\textbf{Aspecto} & \textbf{Modelo AR} & \textbf{MDLM} \\
\midrule
Mecanismo de atención & Atención causal & Atención bidireccional \\
Entrada & Texto limpio $\mathbf{x}$ & Secuencia con ruido $\mathbf{z}_t$ (después de máscaras aleatorias) \\
Muestreo adicional & Ninguno & Se requiere muestrear el paso de tiempo $t$ \\
Función de pérdida & Entropía cruzada & Entropía cruzada ponderada (peso $\frac{\alpha_t'}{1-\alpha_t}$) \\
Posición de cálculo & Todas las posiciones & Solo las posiciones con máscara \\
\bottomrule
\end{tabular}
\end{table}

\begin{knowledgebox}{Concisión del entrenamiento}
Si sabes cómo entrenar un modelo AR, básicamente ya sabes cómo entrenar uno MDLM. La única operación adicional es: (1) Muestrear aleatoriamente un paso de tiempo $t$; (2) Aplicar máscaras aleatorias a la entrada según $t$; (3) Incluir el término de ponderación en la pérdida. Esta concisión es una razón importante por la que MDLM ha sido ampliamente adoptado.
\end{knowledgebox}

\subsection{Detalles del proceso de muestreo}

El proceso completo de muestreo de MDLM (muestreo ancestral, Ancestral Sampling):

\begin{enumerate}
    \item \textbf{Inicialización}: $\mathbf{z}_T = [\text{M}, \text{M}, \ldots, \text{M}]$ (secuencia completamente enmascarada)
    \item \textbf{Denoising de un paso} (de $\mathbf{z}_t$ a $\mathbf{z}_s$):
    \begin{enumerate}
        \item Enviar $\mathbf{z}_t$ a un Transformer bidireccional
        \item El modelo predice la distribución de palabras para todas las posiciones con máscara
        \item Muestrear de la distribución para rellenar todas las posiciones con máscara
        \item Decidir qué predicciones conservar según la probabilidad de desmascaramiento $\frac{\alpha_s - \alpha_t}{1 - \alpha_t}$
        \item Las predicciones no conservadas se vuelven a enmascarar
        \item Los tokens revelados permanecen sin cambios
    \end{enumerate}
    \item Repetir el paso 2 un total de $T$ veces hasta que $\mathbf{z}_0$ esté completamente revelado
\end{enumerate}

\begin{warningbox}{Esto no es un método heurístico}
Este proceso de muestreo no es un método heurístico diseñado por intuición, sino un muestreador ancestral (Ancestral Sampler) derivado estrictamente de deducciones matemáticas. Las deducciones matemáticas requieren que: (1) Los tokens revelados deben permanecer sin cambios; (2) Solo las predicciones en posiciones con máscara se vuelven a enmascarar.
\end{warningbox}

\subsection{Resumen de la sección}

La función de pérdida de entrenamiento de MDLM es una entropía cruzada ponderada, donde el peso está determinado por la probabilidad de desmascaramiento. El límite continuo es la simplificación central de MDLM en comparación con trabajos anteriores. El flujo de entrenamiento es muy similar al de los modelos AR, lo que reduce la barrera para su implementación.

\newpage
\section{Resultados experimentales y aplicaciones}

\subsection{Rendimiento en modelado de lenguaje}

\subsubsection{Conjunto de datos LM1B}

En el conjunto de datos LM1B (longitud de contexto 128), MDLM supera significativamente a todas las líneas base de modelos de difusión discretos previos (incluyendo la línea base más cercana SEDD) y \textbf{acercó la perplexidad de los modelos AR}.

\begin{knowledgebox}{¿Por qué es importante la perplexidad}
La perplexidad es el indicador de evaluación más central en el modelado de lenguaje. La experiencia demuestra que si un modelo logra buenos puntajes de perplexidad en los conjuntos de entrenamiento/validación, su rendimiento en tareas aguas abajo (como programación o razonamiento) también será mejor. Por lo tanto, la perplexidad es una métrica proxy confiable para medir la calidad de los modelos de lenguaje.
\end{knowledgebox}

Cuando la longitud del contexto se aumenta a 1024, MDLM sigue superando a SEDD, pero existe cierta brecha con respecto a los modelos AR. No obstante, dado que las líneas base de modelos de difusión tempranos mostraban una diferencia enorme frente a las líneas base AR, este resultado ya es muy alentador.

\subsubsection{Capacidad de transferencia sin ejemplos}

Después de entrenar el modelo en OpenWebText, se evalúa la verosimilitud sin ejemplos (zero-shot) en 7 conjuntos de datos no vistos: PTB, arXiv, PubMed, entre otros:
\begin{itemize}
    \item MDLM supera consistentemente a SEDD en todos los conjuntos de datos.
    \item En 3 de los 7 conjuntos de datos, el rendimiento de MDLM es \textbf{superior a la línea base AR}.
\end{itemize}

\subsection{Validación a gran escala: LaDa (MDLM de 8B)}

LaDa (Large Language Diffusion Models) es la validación de MDLM a una escala de 8 mil millones de parámetros. Es esencialmente un modelo MDLM de 8B.

\begin{importantbox}{Rendimiento de MDLM a gran escala}
\begin{itemize}
    \item \textbf{Comprensión multidepósito del lenguaje} (MMLU): MDLM \textbf{supera consistentemente la línea base AR} durante todo el proceso de entrenamiento.
    \item \textbf{Tareas de razonamiento}: La curva de expansión del rendimiento de MDLM muestra un crecimiento \textbf{exponencial}, mientras que la línea base AR solo muestra un crecimiento lineal.
    \item \textbf{Preguntas de sentido común y programación}: MDLM es \textbf{emparejado} o \textbf{ligeramente inferior} a la línea base AR.
\end{itemize}
A pesar de que aún existe una brecha en términos de perplexidad, en tareas aguas abajo prácticas, MDLM supera al modelo AR en múltiples aspectos.
\end{importantbox}

\subsection{Más allá del lenguaje: modelado de datos discretos}

\begin{knowledgebox}{Sesgo izquierda-derecha y estructura de datos}
El lenguaje natural tiene un sesgo inherente de izquierda a derecha—esta es la forma en que los humanos hablan y escriben. Por lo tanto, los modelos AR tienen una ventaja natural en tareas de lenguaje. Sin embargo, en áreas como la generación de proteínas, el descubrimiento de fármacos o la generación molecular, los datos \textbf{no presentan sesgo izquierda-derecha}, lo cual es precisamente el escenario donde MDLM destaca.
\end{knowledgebox}

En tareas de generación molecular (artículo GenMol), los modelos basados en MDLM mostraron una ventaja significativa: mediante la regulación de la temperatura de muestreo, se puede obtener una frontera de Pareto entre la \textbf{calidad} y la \textbf{diversidad} de las moléculas. En contraste, la línea base AR solo puede proporcionar un punto fijo de calidad-diversidad e incapaz de realizar este tipo de compensación flexible.

\subsection{Generación controlada}

MDLM es naturalmente adecuado para la generación controlada. La razón es:

\begin{enumerate}
    \item Durante el proceso de eliminación de ruido, el modelo rellena todas las posiciones enmascaradas, produciendo una \textbf{predicción completa} de la muestra final.
    \item De esto se puede obtener información de alto nivel sobre la muestra final.
    \item Utilizando esta información, se pueden guiar muestras con atributos específicos durante el proceso de muestreo.
\end{enumerate}

\begin{warningbox}{Los modelos AR tienen dificultades para realizar muestreo controlado}
Los modelos AR tienen dificultades conocidas en términos de generación controlada. Debido a que, durante el proceso de generación palabra por palabra, el modelo no puede prever la estructura global futura, es muy difícil imponer restricciones globales durante el proceso de generación. El mecanismo de atención bidireccional y la predicción global de eliminación de ruido de MDLM le confieren una ventaja esencial en la generación controlada.
\end{warningbox}

\subsection{Resumen de la sección}

Los resultados experimentales indican: (1) MDLM se acerca al rendimiento de AR a pequeña escala, e incluso supera a AR en tareas aguas abajo a gran escala (especialmente razonamiento); (2) En datos discretos sin sesgo izquierda-derecha (como generación molecular), el rendimiento de MDLM es particularmente destacado; (3) La flexibilidad de muestreo de MDLM lo hace adecuado para aplicaciones como la generación controlada.

\newpage

````
\section{Limitaciones de MDLM: el problema del caché KV}

\subsection{Repaso del caché KV}

En los modelos AR, el \textbf{caché KV} (Key-Value Caching) es la técnica clave para acelerar la inferencia.

\begin{knowledgebox}{Principio de funcionamiento del caché KV}
En un Transformer AR:
\begin{enumerate}
    \item Debido al uso de atención causal, los valores de activación de cada token dependen únicamente del contexto izquierdo
    \item Los valores de activación de los tokens ya generados \textbf{nunca cambian}, porque los tokens futuros no los afectan
    \item Por lo tanto, se puede almacenar en caché (congelar) las claves y valores calculados
    \item Al generar un nuevo token, solo es necesario realizar una propagación hacia adelante para un \textbf{único token}, no para todo el contexto
\end{enumerate}
Esto permite que los modelos AR tengan un costo computacional muy bajo por paso, incluso si requieren miles de pasos de decodificación.
\end{knowledgebox}

\subsection{Por qué MDLM no admite caché KV}

MDLM utiliza atención bidireccional, lo que significa que \textbf{los valores de activación de cada token dependen de todos los demás tokens}.

Considere el siguiente escenario:
\begin{enumerate}
    \item Estado inicial: $[\text{M}, \text{M}, \text{M}, \text{M}]$
    \item Después del primer paso de desruido: $[\text{M}, \text{B}, \text{M}, \text{D}]$ (se revelan las posiciones 2 y 4)
    \item Después del segundo paso de desruido: $[A, \text{B}, \text{M}, \text{D}, \text{M}, \text{F}]$ (se revelan las posiciones 1 y 6)
\end{enumerate}

En el segundo paso, aunque $\text{B}$ y $\text{D}$ no han cambiado, debido a la aparición de $A$ y $F$, bajo atención bidireccional los valores de activación de $\text{B}$ y $\text{D}$ cambiarán. Por lo tanto, \textbf{no se puede congelar el valor de activación en ninguna posición}.

\begin{warningbox}{Cuello de botella de inferencia de MDLM}
Aunque MDLM puede generar texto con muchos menos pasos que la longitud del secuencia (por ejemplo, generar 1000 palabras con 400-500 pasos), cada paso requiere realizar una propagación hacia adelante completa para \textbf{todo el secuencia}. Para secuencias largas, la complejidad $O(n^2)$ de la atención hace que el costo por paso sea mucho mayor que el de un solo paso en un modelo AR (con caché KV). Por lo tanto, en términos de tiempo real de inferencia, MDLM podría ser incluso más lento que un modelo AR.
\end{warningbox}

\subsection{Resumen de la sección}

La ausencia del caché KV es la limitación práctica más grande de MDLM. Aunque la atención bidireccional otorga al modelo una capacidad de percepción global, también provoca que los valores de activación se actualicen globalmente ante cambios en las entradas, lo que obstaculiza el uso del caché KV.

\newpage
\section{SoLM: un modelo de lenguaje basado en la fusión del AR y la difusión}

\subsection{Idea central: ordenamiento + atención causal}

La contribución fundamental de SoLM (Sorting-based Language Model, también conocido como Esoteric Language Model) es introducir soporte para el \textbf{lote} en MDLM. La idea clave es:

\begin{importantbox}{Estrategia de entrenamiento de SoLM}
Reemplazar la atención bidireccional de MDLM por \textbf{atención causal}, lográndolo mediante las siguientes operaciones:
\begin{enumerate}
    \item Realizar un \textbf{ordenamiento/barajado} del secuencia de ruido $\mathbf{z}_t$
    \item Colocar todos los \textbf{tokens limpios a la izquierda} (como contexto conocido)
    \item Colocar todos los \textbf{tokens enmascarados a la derecha} (como objetivos por predecir)
    \item Aplicar \textbf{atención causal} sobre la secuencia ordenada
\end{enumerate}
\end{importantbox}

Con esta disposición:
\begin{itemize}
    \item El primer token enmascarado depende únicamente del contexto limpio
    \item El segundo token enmascarado depende del contexto limpio + el anterior token enmascarado
    \item Y así sucesivamente...
    \item Ningún token enmascarado depende de los tokens enmascarados "futuros" a su derecha
\end{itemize}

\subsection{Proceso de inferencia y caché KV}

El proceso de inferencia de SoLM aprovecha plenamente las capacidades de la caché KV derivadas de la atención causal:

\begin{enumerate}
    \item El muestreador de antepasados decide qué posiciones enmascaradas deben ser denoiseadas (por ejemplo, las posiciones 3 y 2)
    \item Se realiza una \textbf{propagación hacia adelante únicamente para estos tokens enmascarados}, incluyendo sus codificaciones de posición correspondientes
    \item Las activaciones de los tokens ya predichos se \textbf{congelan} (se almacenan en caché)
    \item En el siguiente paso de denoising, solo es necesario procesar los nuevos tokens enmascarados, aprovechando la caché KV guardada
\end{enumerate}

\begin{importantbox}{Ahorro computacional en la inferencia de SoLM}
En comparación con MDLM, que requiere procesar toda la secuencia en cada paso (incluyendo numerosos tokens enmascarados que no transportan información), SoLM \textbf{solo procesa los tokens por denoising} en cada paso. Esto genera ahorros en dos frentes:
\begin{enumerate}
    \item La longitud de la secuencia para la propagación hacia adelante se reduce drásticamente
    \item La caché KV evita el cálculo repetido de las activaciones de los tokens ya procesados
\end{enumerate}
Ambos factores logran conjuntamente una \textbf{aceleración de la inferencia de orden de magnitud}.
\end{importantbox}

\subsection{Interpolación entre AR y difusión}

Otra contribución importante de SoLM es su capacidad para realizar una \textbf{interpolación suave} entre el modo puro AR y el modo puro de difusión:

\begin{itemize}
    \item \textbf{Modo puro AR}: en cada paso solo se denoisa un token (de izquierda a derecha), ofreciendo la mayor calidad pero con la menor velocidad
    \item \textbf{Modo puro de difusión}: en cada paso se denoisear múltiples tokens (en paralelo), logrando la máxima velocidad aunque pueda haber pérdida de calidad
    \item \textbf{Modo intermedio}: permite un equilibrio flexible entre velocidad y calidad según las necesidades
\end{itemize}

\begin{knowledgebox}{Equilibrio flexible entre velocidad y calidad}
Esta capacidad de interpolación significa que el usuario puede seleccionar el punto de operación óptimo según los requisitos del escenario de aplicación específico. Por ejemplo, en servicios en línea sensibles a la latencia se puede optar por una generación paralela más agresiva, mientras que en escenarios con exigencias extremas de calidad se adopta un modo cercano al AR puro.
\end{knowledgebox}

\subsection{Resumen de la sección}

A través de un diseño ingenioso basado en ordenamiento + atención causal, SoLM logra el soporte para caché KV manteniendo la capacidad de generación paralela de los modelos de difusión. Esto permite una mejora de la eficiencia de inferencia de orden de magnitud, haciendo que los modelos de lenguaje basados en difusión finalmente alcancen velocidades de inferencia prácticas para su uso real.

\newpage
\section{Comparación de métodos y análisis de rendimiento}

\subsection{MDLM frente a BD3LM y SoLM}

BD3LM (Modelo de lenguaje de difusión por bloques) es otro método de modelo de lenguaje de difusión que admite caché KV. Su enfoque consiste en dividir secuencias largas en bloques de tamaño fijo (por ejemplo, 100 tokens por bloque), realizar el proceso de desruido mediante difusión dentro de cada bloque y avanzar entre bloques utilizando un modo autorregresivo.

\begin{table}[H]
\centering
\caption{Comparación de tres métodos de modelo de lenguaje de difusión}
\begin{tabular}{@{}lccc@{}}
\toprule
\textbf{Característica} & \textbf{MDLM} & \textbf{BD3LM} & \textbf{SoLM} \\
\midrule
Tipo de atención & Bidireccional & Bidireccional dentro del bloque + causal entre bloques & Causal (tras ordenar) \\
Caché KV & No admite admitida parcialmente (solo entre bloques) & Admitida completamente \\
Generación paralela & Paralelismo global & Paralelismo dentro del bloque & Paralelismo global \\
Velocidad de inferencia & Lenta & Moderada & Rápida \\
Calidad de generación & Alta & Media-alta & Alta \\
\bottomrule
\end{tabular}
\end{table}

\begin{warningbox}{Degradación de calidad de BD3LM en muestreo con pocos pasos}
BD3LM mantiene una buena calidad al aumentar el número de pasos de muestreo, pero su calidad se degrada significativamente cuando se reduce drásticamente el número de pasos de muestreo (es decir, generar más palabras en paralelo por paso), e incluso puede ser inferior a la de MDLM. Esto se debe a que, tras reducir los pasos de difusión dentro del bloque, el efecto de desruido es insuficiente. En comparación, SoLM mantiene una buena calidad incluso con muestreo reducido.
\end{warningbox}

\subsection{Curva de compensación entre velocidad y calidad}

Sobre la frontera de Pareto entre perplexidad y tiempo de muestreo:

\begin{itemize}
    \item \textbf{MDLM}: Perplexidad óptima, pero velocidad de inferencia más lenta (la curva se sitúa en la parte inferior derecha)
    \item \textbf{BD3LM}: Es más rápido que MDLM con un alto número de pasos de muestreo, pero la calidad disminuye drásticamente tras reducir los pasos
    \item \textbf{SoLM}: La frontera de Pareto es \textbf{estrictamente superior} a la de MDLM y BD3LM; es decir, ofrece mayor calidad para cualquier velocidad dada o mayor velocidad para cualquier calidad dada
\end{itemize}

En pruebas basadas en 110 millones de parámetros y un contexto de longitud 1024 tokens, el tiempo de muestreo de SoLM logró una reducción significativa en comparación con MDLM.

\subsection{Resumen de la sección}

SoLM es estrictamente superior a MDLM y BD3LM en términos de compensación entre velocidad y calidad, lo cual se debe principalmente a su soporte completo para caché KV. Aunque el diseño por bloques de BD3LM también admite parcialmente la caché KV, sigue estando limitado por los cuellos de botella de la atención bidireccional dentro del bloque.

\newpage

\section{Síntesis y ampliación}

\subsection{Repaso de los puntos clave}

\begin{enumerate}
    \item Los \textbf{modelos de lenguaje por difusión} rompen la restricción de generación izquierda a derecha de los modelos AR, permitiendo generar texto de manera paralela.
    \item El \textbf{MDLM} simplifica el entrenamiento de la difusión discreta haciéndolo tan sencillo como en los modelos AR mediante un ELBO continuo y conciso, logrando un rendimiento a gran escala que puede igualar o incluso superar al de los modelos AR.
    \item La \textbf{memoria KV} es el mayor cuello de botella para la implementación real del MDLM: la atención bidireccional impide congelar los valores de activación.
    \item El \textbf{SoLM}, mediante un diseño que combina ordenamiento y atención causal, mantiene la capacidad de generación paralela mientras soporta la memoria KV, logrando una aceleración en el razonamiento de varios órdenes de magnitud.
    \item En tareas con datos discretos sin sesgo izquierda-derecha (generación de proteínas, moléculas), las ventajas de los modelos de lenguaje por difusión son especialmente notables.
\end{enumerate}

\subsection{Estado actual de la aplicación industrial}

Hasta el momento del grabado de esta conferencia (otoño de 2025), el marco MDLM ha sido ampliamente adoptado en la industria:

\begin{itemize}
    \item \textbf{ByteDance Seed Diffusion}: Construido sobre el marco MDLM.
    \item \textbf{Google Gemini Diffusion}: Utiliza métodos de difusión para la generación de texto.
    \item \textbf{Inception Labs Mercury Coder}: Modelo de difusión orientado a la generación de código.
    \item \textbf{NVIDIA}: Emplea MDLM para la generación de moléculas (ámbito del descubrimiento de fármacos).
\end{itemize}

\subsection{Lecturas adicionales}

\begin{itemize}
    \item Sahoo et al., ``Simple and Effective Masked Diffusion Language Models'' (MDLM), NeurIPS 2024.
    \item Sahoo et al., ``Esoteric Language Models'' (SoLM), 2025.
    \item ``Difusión dualidad'': algoritmo de destilación que genera 1000 palabras en 10 pasos.
    \item LaDa: Large Language Diffusion Models (validación en escala de 8B para MDLM).
    \item BD3LM: Block Diffusion Language Models.
    \item GenMol: Generación de moléculas basada en MDLM.
    \item Página del curso KAIST CS492D: \url{https://mhsung.github.io/kaist-cs492d-fall-2025/}
\end{itemize}

%% --- Fin del contenido --- %%

\end{document}
